\documentclass{article}
\usepackage{amsmath,amssymb}
\usepackage{xurl}
\usepackage[hidelinks]{hyperref}
% Shared commands for notes in docs/paper-gaps/.

\DeclareUnicodeCharacter{2080}{\ifmmode{}_{0}\else\textsubscript{0}\fi}
\DeclareUnicodeCharacter{2081}{\ifmmode{}_{1}\else\textsubscript{1}\fi}
\DeclareUnicodeCharacter{2082}{\ifmmode{}_{2}\else\textsubscript{2}\fi}
\DeclareUnicodeCharacter{2099}{\ifmmode{}_{n}\else\textsubscript{n}\fi}

\newcommand{\Lean}{\textsc{Lean}}
\ExplSyntaxOn
\NewDocumentCommand{\leanid}{m}{
  \ifmmode
    \textsf{\detokenize{#1}}
  \else
    \group_begin:
    \urlstyle{sf}
    \tl_set:Nn \l_tmpa_tl {#1}
    \tl_replace_all:Nnn \l_tmpa_tl {\_} {_}
    \exp_args:NV \path \l_tmpa_tl
    \group_end:
  \fi
}
\ExplSyntaxOff
\providecommand{\tr}{\operatorname{tr}}
\newcommand{\tnleanIssueBase}{https://github.com/LionSR/TNLean/issues/}
\newcommand{\tnleanPrBase}{https://github.com/LionSR/TNLean/pull/}
\newcommand{\tnleanDocBase}{https://sirui-lu.com/TNLean/docs/}
\newcommand{\ghissue}[1]{\href{\tnleanIssueBase#1}{GitHub issue \##1}}
\newcommand{\ghpr}[1]{\href{\tnleanPrBase#1}{PR \##1}}
\newcommand{\doclink}[1]{\href{\tnleanDocBase#1.html}{\path{#1}}}

% Dirac notation and the identity operator, as in blueprint/src/macros/common.tex.
% \providecommand keeps note-local definitions authoritative.
\usepackage{dsfont}
\providecommand{\ket}[1]{|#1\rangle}
\providecommand{\bra}[1]{\langle #1|}
\providecommand{\braket}[2]{\langle #1 | #2 \rangle}
\providecommand{\ketbra}[2]{|#1\rangle\!\langle #2|}
\providecommand{\Id}{\mathds{1}}
\providecommand{\one}{\mathds{1}}
\providecommand{\C}{\mathbb{C}}
\providecommand{\MN}[1]{M_{#1}(\C)}
\providecommand{\MD}{M_D(\C)}

% External referencing for paper sources.  The build helper writes sanitized
% auxiliary files under build/paper-aux/ and excludes duplicated source labels.
\usepackage{xr}

% Import a generated paper auxiliary file with a caller-supplied label prefix.
% The candidate paths support builds from docs/paper-gaps/, the repository root,
% and build/paper-aux/ itself.
\newcommand{\importpaperaux}[2]{%
  \IfFileExists{../../build/paper-aux/#2.aux}{%
    \externaldocument[#1]{../../build/paper-aux/#2}%
  }{%
    \IfFileExists{build/paper-aux/#2.aux}{%
      \externaldocument[#1]{build/paper-aux/#2}%
    }{%
      \IfFileExists{#2.aux}{%
        \externaldocument[#1]{#2}%
      }{}%
    }%
  }%
}

% General external reference macros
\newcommand{\paperref}[2]{\ref{#1-#2}}
\newcommand{\papereqref}[2]{(\ref{#1-#2})}

% CPSV16 source (1606.00608 / MPDO-22-12-17-2.tex) external document and shortcuts
\importpaperaux{cpsv16-}{1606.00608/MPDO-22-12-17-2-tnlean}
\newcommand{\cpsvref}[1]{\paperref{cpsv16}{#1}}
\newcommand{\cpsveqref}[1]{\papereqref{cpsv16}{#1}}

% Verdict marker: \gapnote{<kind>}{<status>}. Kind is the policy
% classification (clarification, local-correction, scope-restriction,
% unfaithful, false-source, open-gap); status is open, wip (elimination
% underway), resolved, or historical. Machine-read by the site index;
% prints a verdict line.
\newcommand{\gapnote}[2]{\par\noindent\textbf{Verdict:} #1 (#2).\par}

\title{The RVB Tensor: the Singlet Matrix on the Vacuum}
\author{}
\date{2026-09-26}
\begin{document}
\maketitle
\gapnote{local-correction}{resolved}

\section{What the source prints}
Appendix~A of \cite{Cirac2021Matrix}, under ``Two dimensions: PEPS'', writes
the nearest-neighbour resonating valence bond state as a PEPS of bond
dimension three. Each site carries the projector
\begin{align}
    P=\ket{0}\bigl[(0222|+(2022|+\dots\bigr]
    +\ket{1}\bigl[(1222|+(2122|+\dots\bigr],
    \label{eq:rmp_peps_rvb_projector}
\end{align}
whose nonzero entries have exactly one virtual index in a spin state $0,1$
and the other three in the vacuum state $2$. The projector is combined with
$\pm Y$ on the links, and on the square lattice with a translation-invariant
orientation of the singlets the tensor is
$A=P(\mathbb 1\otimes\mathbb 1\otimes Y\otimes Y)$, where
$Y=\left(\begin{smallmatrix}0&-1\\1&0\end{smallmatrix}\right)$ is the
$2\times 2$ singlet matrix introduced for the AKLT
state.\footnote{Source traceability:
    \path{Papers/2011.12127/TN-Review-main.tex}, lines 2434 and
    2440--2448.}

\section{The degenerate reading}
The bond has dimension three, while $Y$ is printed as a $2\times 2$ matrix.
If $Y$ is extended to the bond by zero on the vacuum, the down and left legs
of $A$ can never carry the vacuum index: every column of
$\mathbb 1\otimes\mathbb 1\otimes Y\otimes Y$ with a vacuum down or left
index vanishes, and every other column has at least two spin indices, which
\eqref{eq:rmp_peps_rvb_projector} annihilates. The tensor $A$ is then zero.
The source plainly does not intend this reading, since it describes the
state as the superposition of all nearest-neighbour singlet coverings.

\section{The correction}
The singlet matrix acts on the bond as $Y\oplus 1$: it is $Y$ on the spin
states $0,1$ and the identity on the vacuum $2$, so the state on each bond is
$\ket{01}-\ket{10}+\ket{22}$. An edge with vacuum at both ends carries no
singlet, and an edge with spin labels at both ends carries a singlet. With
this reading, on a torus of width and height at least three, the coefficient
of a configuration $\sigma$ is the sum over nearest-neighbour dimer coverings
$c$ of the product of the singlet coefficients $Y_{\sigma_{v+e},\sigma_v}$
over the edges $(v,v+e)$ of $c$, with $e$ the unit step to the right or
up.\footnote{Results: \leanid{TNLean.PEPS.rvbSingletY},
    \leanid{TNLean.PEPS.stateCoeff_rvbPEPS}.}

\section{Elimination plan}
None needed beyond this record: the extension by the identity on the vacuum
is part of the definition of the tensor. Any other nonzero entry $c$ on the
vacuum gives the same state up to normalization, since each edge not in a
covering carries the vacuum on exactly one down or left leg and every
covering of a torus of $N$ sites leaves $3N/2$ edges uncovered, so each
covering acquires the same factor $c^{3N/2}$.

\bibliographystyle{alpha}
\bibliography{references}
\end{document}
